
\subsection{TPU slices and device placement}
\label{sec:bg-tpu}

A single Cloud TPU v5e chip, the generation used throughout this
work, provides 197\,bf16 TFLOPs and 16\,GB of HBM with a memory bandwidth of 800\,GiB/s
\citep{googletpuv5e}. Multiple chips on a slice are connected by a
dedicated Inter-Chip Interconnect (ICI) fabric arranged as a 2D torus,
independent of the host network, allowing up to 256 chips per pod to
exchange activations and gradients without going through the datacenter
network \citep{googletpuv5e}. Because v5e has no dedicated
peer-reviewed architecture paper, we cite the vendor documentation for
per-chip specifications and the two Jouppi ISCA papers
\citep{jouppi2017tpu, jouppi2023tpuv4} for the systolic-array design and
its evolution across TPU generations.

Software sees each chip as an explicit JAX device. Nothing about the
torus interconnect is used unless the program says so. A slice with
$n$ chips visible does not automatically distribute an $n$-chip job
across them: how work is mapped onto devices is decided by the JAX/XLA
compilation layer, or explicitly by the programmer, not by the
hardware. Section~\ref{sec:parallelism} measures what that layer does
with AlphaFold2 on this slice.

\subsection{The JAX/XLA compilation model}
\label{sec:bg-jax}

JAX traces Python functions into an intermediate representation
(a \texttt{jaxpr}) the first time they are called with a new
combination of argument shapes, then lowers that representation to
XLA HLO and compiles it to a binary for the target device. The
profile labels the function that performs this trace-and-compile step
\texttt{cache\_miss}, after the fallback path taken when JAX's
fast-path jit cache does not hold an entry for the call; for a
Haiku-based model such as AlphaFold2, this trace walks through
Haiku's \texttt{apply\_fn}, lowering the full Evoformer computation
graph into XLA's representation before a single matrix multiply has
executed. Compilation is keyed on input shapes and dtypes, together
with any static arguments: a function called again with the same
signature reuses the compiled binary, while a change to it (a different
sequence length, or a switch between \texttt{float32} and
\texttt{bfloat16}) triggers a fresh trace and compile. A workload's
\emph{first} call at a given input shape is therefore different work
from every call after it, and the timing methodology has to separate
the two (Section~\ref{sec:methodology}). A persistent compilation
cache can carry compiled binaries across process restarts, avoiding
the XLA-compile step on a shape that has already been compiled once,
but a fresh process still has to re-trace the Python program down to a
jaxpr before it can consult that cache, so a persistent cache reduces
cold-start cost without eliminating it.

JAX's own transformations for expressing
parallelism, \texttt{vmap} for vectorizing a function over a batch axis
and \texttt{pmap} for mapping a function across physical devices with
an explicit per-device argument axis, are different operations even
though both add a leading axis to a function's inputs: \texttt{vmap}
changes what a single device computes, \texttt{pmap} changes which
device computes it. XLA's automatic partitioner, in turn, can shard a
computation across devices
without a user manually writing per-device code, but it does so by
propagating sharding decisions outward from explicit annotations,
principally \texttt{PartitionSpec} constraints placed on inputs,
outputs or intermediates. (JAX also offers \texttt{shard\_map}, the
manual complement to automatic partitioning: there the programmer
writes the per-device code and the collectives.) A computation carrying no annotations gives the partitioner
nothing to propagate from, and its documented fallback is to replicate
the computation rather than split it.

\subsection{AlphaFold2's computational structure}
\label{sec:background-af2}

AlphaFold2 predicts a protein's three-dimensional structure from its
amino-acid sequence by iterating two representations against each
other: a multiple sequence alignment (MSA) representation, one row per
homologous sequence, and a pair representation, one entry per pair of
residues, encoding the model's evolving belief about which residues
are close in space. The Evoformer block updates both representations
through attention operations (including triangular self-attention and
triangular multiplicative updates on the pair representation, whose
cost scales with the number of residues beyond simple quadratic
attention) before a structure module converts the final single and pair
representations into three-dimensional coordinates
\citep{jumper2021alphafold}. The whole Evoformer-plus-structure-module
pipeline is then \emph{recycled}: its own output is fed back as input
for a further pass, refining the prediction over a small, fixed number
of iterations \citep{jumper2021alphafold}. AlphaFold2's public
implementation is written in JAX and Haiku, which is what makes it
usable, unmodified, as a systems benchmark across JAX-supported
backends (Section~\ref{sec:methodology}): a real, widely used model
with an attention-heavy graph and small sequential control flow.
Its diffusion-based successor,
AlphaFold3, replaces the Evoformer/structure-module pair with a
different architecture, and its released entry point accepts a
different set of backends (Section~\ref{sec:alphafold3}). We report
those as two separate facts: nothing we measured connects the one to
the other.

One further structural detail matters for multi-chip execution.
AlphaFold2 can run an \emph{ensemble}: several stochastic variations of
the same input are processed and their representations averaged. The
released configuration uses a single member for evaluation, as do our
runs. When more members are requested, the public implementation does
not compute them along an array axis. It iterates over them with a
sequential \texttt{hk.while\_loop} inside \texttt{AlphaFoldIteration},
and recycling is implemented as a \texttt{while\_loop} in the same way.
Sequential control flow of this kind offers no tensor dimension to
split across devices, so the ensemble cannot be parallelized by
sharding its arrays; spreading members across chips requires mapping
them explicitly, which is what we do in Section~\ref{sec:parallelism}.
This is a separate matter from the partitioner behavior we observe on a
single-member forward pass, for which the most plausible explanation
is the general mechanism described in Section~\ref{sec:bg-jax}: the
model carries no sharding annotations for the partitioner to propagate.
Section~\ref{sec:parallelism} states what evidence we have for that
and what we lack.
